\section*{الفصل \ref{ch:b1:mammal-organization} --- التنظيم الوظيفي عند الثدييات}

\begin{solution}{exo:b1:mammal-organization:1}
عند الفقاريات: التناظر الجانبي، \omterm{def:b1:mammal-organization:bodyplan}{والتمركز الرأسي}، والمحور الفقري
الظهري \omterm{def:b1:organism-environment:organism}{المحيط} بالنخاع الشوكي، \omterm{def:b1:mammal-organization:bodyplan}{والجوف العام} البطني بأحشائه، والأطراف
الأربعة. وعند الثدييات: \omterm{def:b1:mammal-organization:bodyplan}{الحجاب الحاجز} يقسم \omterm{def:b1:mammal-organization:bodyplan}{الجوف العام} إلى صدر وبطن،
والشعر، واللبن، وحرارة الجسم المنظَّمة.
\end{solution}

\begin{solution}{exo:b1:mammal-organization:2}
التغذية: الهضمي (المعدة)، والتنفسي (الرئة)، والدوراني (القلب)،
والإطراحي (الكلية). العلاقة: العصبي (الدماغ)، \omterm{def:b1:mammal-organization:organ}{وأعضاء} الحس (العين)،
والصماوي (الدرقية)، والهيكل والعضلات (عظم الفخذ، العضلة ذات
الرأسين)، والجلد. التكاثر: التناسلي (المبيض، الخصية)، والغدد اللبنية.
\end{solution}

\begin{solution}{exo:b1:mammal-organization:3}
الماء $0.6\times 60 = \qty{36}{L}$؛ وداخل الخلوي \qty{24}{L}؛ وخارج
الخلوي \qty{12}{L}، منه الخلالي \qty{9}{L} والبلازما \qty{3}{L}.
\end{solution}

\begin{solution}{exo:b1:mammal-organization:4}
$1.1/5 = 22\%$. و\qty{1.1}{L/min} هي \qty{1580}{L} في اليوم، أي حجم
الدم 290 مرة.
\end{solution}

\begin{solution}{exo:b1:mammal-organization:5}
الإينولين الباقي \qty{1.7}{g}؛ فيكون $V = 1.7/0.12 = \qty{14}{L}$
خارج الخلوي؛ وداخل الخلوي $42 - 14 = \qty{28}{L}$.
\end{solution}

\begin{solution}{exo:b1:mammal-organization:6}
جدار الشعرية يمرر الماء والشوارد الصغيرة بحرية فتتوازن، لكنه يحبس
البروتينات. وبروتينات البلازما تجذب الماء إلى الأوعية تناضحيًا؛ ولولاها
لرشح الماء إلى الحيز الخلالي فتورم (وذمة)، ولهبط حجم البلازما.
\end{solution}

\begin{solution}{exo:b1:mammal-organization:7}
على التوازي ينال كل \omterm{def:b1:mammal-organization:organ}{عضو} دمًا بكامل أكسجينه الشرياني وضغطه، ويمكن
تنظيم جريانه على حدة. أما على التسلسل فسينال الدماغ دمًا استنزفته
العضلات، بضغط أدنى، وسيهبط إمداده في أثناء الجهد بالذات.
\end{solution}

\begin{solution}{exo:b1:mammal-organization:8}
$hS\Delta T = 80 - 20 = \qty{60}{W}$؛ فيكون
$h = 60/(1.8\times 13) =
\qty{2.6}{W.m^{-2}.K^{-1}}$. وعند هواء
\qty{33}{\celsius} يزول الفقد غير التبخري؛ فلا بد للجلد أن يسخن فوق
الهواء (توسع الأوعية يرفعه نحو \qty{36}{\celsius}) وأن يحمل التعرق
الباقي.
\end{solution}

\begin{solution}{exo:b1:mammal-organization:9}
الحرارة الواجب طرحها $800 - 200 - 150 = \qty{450}{W}$، أي
\qty{1620}{kJ/h}، أي \qty{675}{g} من العرق في الساعة. ودون تعرق:
$450/(70\times 3500) =
\qty{1.8e-3}{K/s}$، أي نحو \qty{0.11}{K/min}:
أي \qty{3}{K} في نصف ساعة.
\end{solution}

\begin{solution}{exo:b1:mammal-organization:10}
الأجسام الصغيرة تفقد الحرارة أسرع لكل غرام (لكبر $S/V$)، وتحتاج إلى
إنتاج مستمر عالٍ موثوق: فالشحم البني ينتج الحرارة كيميائيًا بمعدل
ثابت دون حركة، بينما يقتضي الارتعاش عضلات متطورة (تفتقر إليها
الوليدة)، ويعوق الحركة والتغذية، وهو متقطع. كما أن الشحم البني يقع
قرب الأوعية الكبيرة فيدفئ الدم مباشرة.
\end{solution}

\begin{solution}{exo:b1:mammal-organization:11}
الوطاء المسخَّن يقرأ «حرّ زائد»: فيلهث الكلب ويوسع أوعية جلده رغم
البرد، وتهبط حرارة لبه بدرجة أو درجتين في غضون الساعة. ومن غير وطاء
لا \omterm{def:b1:organism-environment:homeostasis}{قيمة مرجعية}: فلا يرتعش الكلب في البرد ولا يلهث في الحر، وتنزاح
حرارته نحو حرارة الغرفة في الحالتين.
\end{solution}

\begin{solution}{exo:b1:mammal-organization:12}
\omterm{def:b1:mammal-organization:compartments}{السائل خارج الخلوي} (\omterm{def:b1:mammal-organization:compartments}{السائل الخلالي} والبلازما)، وهو شبه ثابت التركيب
غني بالصوديوم ثابت الحرارة والحجم، يغمر كل خلية كما غمر ماء البحر
أول الحيوانات. وثباته تصونه حلقات تغذية راجعة سالبة (الكلية للحجم
والشوارد، والرئة للغازات، والكبد والبنكرياس للغلوكوز، والوطاء
للحرارة). و«شواطئه» هي سطوح التبادل --- المعي، والرئة، والكلية،
والجلد --- حيث يلاقي الخارج، والدوران هو التيار الذي يحركه.
\end{solution}

\begin{solution}{pb:b1:mammal-organization:1}
\textbf{1.} $B = 3.4\times 2^{0.75} = \qty{5.7}{W}$.
\textbf{2.} في الجحر: $5.7\times 12\times 3600 = \qty{246}{kJ}$؛ وفي
البحث عن الغذاء: $17.1\times 6\times 3600 = \qty{369}{kJ}$؛ وفي
الخارج: $11.4\times 6\times 3600 = \qty{246}{kJ}$.
\textbf{3.} \qty{861}{kJ}، أي $861/493 = 1.75$ مرة المعدل القاعدي
($B$ على مدى \qty{24}{h} $= \qty{493}{kJ}$).
\textbf{4.} $0.25\times 18\times 0.65 = \qty{2.9}{kJ/g}$.
\textbf{5.} $861/2.9 = \qty{295}{g}$ من العشب الطري، أي \qty{74}{g}
من المادة الجافة.
\textbf{6.} $295/2000 = 15\%$ من كتلته.
\textbf{7.} $861/20 = \qty{43}{L}$ من $\mathrm{O_2}$؛ أي
\qty{30}{mL/min}.
\textbf{8.} $0.9\times 43 = \qty{39}{L}$ من $\mathrm{CO_2}$، أي $39/22.4\times 44 = \qty{76}{g}$.
\textbf{9.} $0.75\times 295 = \qty{221}{g}$.
\textbf{10.} المادة الجافة المهضومة $0.65\times 74 = \qty{48}{g}$؛
والماء $0.6\times 48 = \qty{29}{g}$.
\textbf{11.} غير المهضومة \qty{26}{g}؛ والبراز
$26/0.6 = \qty{43}{g}$، يحمل \qty{17}{g} من الماء.
\textbf{12.} $(44 - 5)\times 0.86 = \qty{34}{g}$.
\textbf{13.} الوارد: $221 + 29 = \qty{250}{g}$. والصادر:
$17 + 130 + 34 + 20
= \qty{201}{g}$. والفائض \qty{49}{g}.
\textbf{14.} لا: فالعشب يمده بأكثر مما يفقد؛ ويخرج الفائض بولًا
زائدًا (نحو \qty{180}{mL} في المجموع).
\textbf{15.} المادة الجافة نفسها \qty{74}{g} في $\qty{148}{g}$ من
العشب: فالماء في الغذاء \qty{74}{g}؛ والوارد \qty{103}{g} مقابل صادر
\qty{201}{g}: أي عجز \qty{100}{g} في اليوم، عليه أن يشربه أو يفقده
من جسمه.
\textbf{16.} الماء $0.6\times 2.0 = \qty{1.2}{L}$؛ وخارج الخلوي
\qty{0.4}{L}؛ والبلازما \qty{0.1}{L}.
\textbf{17.} $1.5/15 = \qty{0.10}{L}$: وهو متفق مع التقدير.
\textbf{18.} $0.10/0.60 = \qty{0.167}{L}$، أي \qty{8}{\%} من كتلة
الجسم.
\textbf{19.} $5\times(2/70)^{0.75} = \qty{0.35}{L/min}$؛ وزمن الدورة
$0.167/0.35 = \qty{0.48}{min}$، أي نحو \qty{29}{s}.
\textbf{20.} $350/200 = \qty{1.75}{mL}$ لكل نبضة.
\textbf{21.} $0.1\times 2^{2/3} = \qty{0.16}{m^2}$.
\textbf{22.} $2.5\times 0.16\times 34 = \qty{13.6}{W}$، مقابل
$2B = \qty{11.4}{W}$ المفترضة: فعلى الأرنب أن ينتج أكثر قليلًا، أو
يخفض الفقد بالوضعية والمأوى.
\textbf{23.} مفتوحتان: $10\times 0.01\times 34 = \qty{3.4}{W}$، أي
ربع فقد الفرو زيادةً؛ ومغلقتان:
$10\times 0.01\times 3 =
\qty{0.3}{W}$. فهو يغلقهما.
\textbf{24.} الفرو: $2.5\times 0.16\times 9 = \qty{3.6}{W}$؛
والأذنان مفتوحتان: $10\times 0.01\times 9 = \qty{0.9}{W}$؛
والمجموع \qty{4.5}{W} مقابل \qty{5.7}{W} منتجة. فلا بد للباقي
\qty{1.2}{W} أن يخرج بالتبخر: فيلهث الأرنب (\qty{1.2}{W} هي
\qty{1.8}{g} من الماء في الساعة)، وينبطح على أرض باردة، ويلزم الظل.
\textbf{25.} الإنفاق الطاقي اليومي \qty{0.86}{MJ}، أي $1.75$ مرة
معدله القاعدي؛ ودوران الماء \qty{250}{g}، أي \qty{12.5}{\%} من
كتلته.
\end{solution}
